\documentclass[11pt]{article}
\title{On the Closure of Compositional Hierarchies in Digital Symmetries}
\author{A rigorous researcher-engineer bridging abstract mathematics and concrete digital systems.}
\usepackage[margin=1in]{geometry}
\usepackage{graphicx}
\usepackage{amsmath,amssymb,mathtools}
\usepackage{hyperref}
\graphicspath{{media/}{media/diagrams/}{images/}{artwork/}}

\begin{document}

\section*{Glossary \& Symbol Index}

\begin{description}
  \item[Closure] A property of a class of objects or operations being stable under the relevant composition rules used throughout the book. In the main text, closure is developed in the context of clones, operads, and compositional hierarchies; see the foundational chapters on closure theory and modular composition for formal statements and proofs.

  \item[Clone] A set of Boolean functions closed under composition and containing the projection maps. Clones provide the algebraic language for describing which digital operations can be built from a chosen basis; see the chapter on Post's lattice for the classification results.

  \item[Compositional hierarchy] A layered organization of systems in which higher-level components are assembled from lower-level ones by explicit composition rules. The book uses this term to connect algebraic structure, circuit design, and modular refinement.

  \item[Contract] A specification of assumptions and guarantees for a module or component. Contracts are used to state interface behavior and to support refinement and verification in the modularity chapters.

  \item[Digital symmetry] A transformation of a digital object that preserves the structure of interest, such as a Boolean function, circuit, or dataflow graph. Symmetries may include permutations of inputs or wires, dualities, and other structure-preserving relabelings.

  \item[Equivariance] The property that a map commutes with a symmetry action. Equivariant constructions appear in the chapters on symmetry groups, learning on \(\mathbb{B}^n\), and parallel dataflow.

  \item[Feedback] A composition pattern in which outputs are routed back into inputs. The book treats feedback through sequential composition, traced categories, and stateful machines.

  \item[Functional completeness] The ability of a set of Boolean primitives to generate every Boolean function under composition. The NAND-based case study and the discussion of Post's lattice use this notion as a central benchmark.

  \item[Hierarchy] A stratified organization of components, abstractions, or proofs in which each level is built from the previous one. In this book, hierarchy is always tied to explicit composition and closure conditions rather than informal layering.

  \item[Interface] The boundary description of a component, including its inputs, outputs, and typing or arity constraints. Interfaces are the primary objects of comparison in the modular and contract-based chapters.

  \item[Invariance] The property that a quantity, relation, or behavior remains unchanged under a specified transformation. Invariance is used throughout the book to compare circuits, functions, and metrics across symmetries.

  \item[Module] A component with a declared interface that can be composed with other components. Modules are the basic units of hierarchical design in the book's construction-oriented chapters.

  \item[Operad] An algebraic structure for composing multi-input operations according to arity and substitution rules. Operads provide one of the book's main formalisms for compositional digital structure.

  \item[PROP] A symmetric monoidal category whose objects are natural numbers and whose morphisms represent operations with multiple inputs and outputs. PROPs are used to model circuits, wiring diagrams, and related compositional systems.

  \item[Refinement] A relation between specifications or implementations in which one artifact satisfies a stronger or more concrete version of another. Refinement is used to connect abstract contracts to executable builds.

  \item[Reversible computation] Computation in which the mapping between inputs and outputs is bijective, so that information is not discarded. The reversible and quantum-flavored chapters use this notion to connect digital symmetry with physical constraints.

  \item[Symmetry group] The group of all structure-preserving transformations of a given object. For Boolean functions and circuits, this often includes permutations and input/output relabelings that preserve behavior.

  \item[Trace] An operation that closes a feedback loop in a monoidal setting. Traces formalize sequential composition with feedback and are used to reason about stateful systems.

  \item[\(\mathbb{B}^n\)] The \(n\)-dimensional Boolean cube, i.e., the set of all length-\(n\) bit vectors. It is the ambient space for Boolean functions, symmetry actions, and learning models in the book.
\end{description}

\medskip

\noindent\textbf{Symbol index.} The table below provides a bidirectional map from symbols to their primary locations in the book. Page numbers are intentionally left as placeholders until the final pagination pass; chapter and section references are the stable navigation keys.

\begin{center}
\begin{tabular}{@{}lll@{}}
\hline
Symbol & Meaning & Primary location \\
\hline
\(\mathbb{B}^n\) & Boolean \(n\)-cube & Chapter 14, \emph{Learning on \(\mathbb{B}^n\) and Equivariant Architectures} \\
\(\circ\) & Composition of operations or maps & Front matter: Notation \& Conventions; Chapters 2--7 \\
\(\mathrm{Bool}\) & The Boolean domain / Boolean algebra context & Chapter 3, \emph{Closure \& Clones: Post's Lattice for \(\mathrm{Bool}\)} \\
\(\mathrm{Clone}\) & Composition-closed set of Boolean operations & Chapter 3 \\
\(\mathrm{NAND}\) & Sheffer primitive used for functional completeness & Chapter 7, \emph{Case Study---From NAND to XOR to Half-Adder} \\
\(\mathrm{XOR}\) & Exclusive-or gate / parity operation & Chapters 7, 8, and 14 \\
\(\mathrm{AND}\) & Conjunction gate & Chapters 3 and 7 \\
\(\mathrm{OR}\) & Disjunction gate & Chapters 3 and 7 \\
\(\mathrm{NOT}\) & Negation gate & Chapters 3 and 7 \\
\(\otimes\) & Monoidal tensor / parallel composition & Chapter 2, \emph{Operads, PROPs, and Symmetric Monoidal Foundations} \\
\(\mathrm{Tr}\) & Trace / feedback operator & Chapter 5, \emph{Sequential Composition, Feedback, and Traces} \\
\(\eta\) & Unit / identity-like structure, context dependent & Chapters 2, 4, and 6 \\
\(\epsilon\) & Counit / evaluation-like structure, context dependent & Chapters 2, 4, and 6 \\
\(f\) & Generic function or morphism & Throughout the book; see Notation \& Conventions \\
\(g\) & Generic function or morphism & Throughout the book; see Notation \& Conventions \\
\(\sim\) & Equivalence relation / symmetry equivalence, context dependent & Chapters 8 and 15 \\
\(\cong\) & Isomorphism / structural equivalence & Chapters 2, 4, 12, and 15 \\
\(\vdash\) & Derivability / proof judgment & Chapter 6, \emph{Constructive Logic: Heyting Semantics for Digital Composition} \\
\hline
\end{tabular}
\end{center}

\medskip

\noindent\textbf{Location key.} Chapter references are given by title so that the index remains stable even if pagination changes. When a symbol appears in multiple chapters, the listed location indicates the first or most central discussion; later uses should be read in light of that primary definition.

\medskip

\noindent\textbf{Reading note.} For formal definitions, proofs, and build artifacts, follow the cross-references in the main text rather than relying on this glossary alone. This section is intended as a navigational aid and a plain-language reminder of the book's recurring symbols and terms.

\medskip

\noindent\textbf{How to use this index.} The glossary is organized for two-way lookup: readers can move from a term to its plain-language meaning, then from a symbol to the chapter where it is first defined or used centrally. When a symbol is reused across several chapters, the index points to the most relevant entry point rather than attempting to enumerate every occurrence.

\medskip

\noindent\textbf{Scope note.} This section is intentionally descriptive rather than exhaustive. It records the recurring vocabulary needed to navigate the book's proofs, constructions, and build artifacts, but it does not replace the formal definitions given in the chapters themselves.


\end{document}
